\chapter{Einleitung}
\label{ch:introduction}

\section{Motivation und Problemstellung}
\label{sec:motivation}

Die Wiederverwendung von Altholz in Tragwerken bietet ein hohes Potential fuer
Ressourceneffizienz und Kreislaufwirtschaft im Bauwesen. Nach Jahrzehnten des
Einsatzes unter Belastung bleibt eine zentrale Frage ungeklaert: Unterscheidet
sich die Restfestigkeit zwischen verschiedenen Bereichen (Zonen) eines
gebrauchten Holzbalkens signifikant?

Ein Holzbalken unter gleichmaessiger Belastung erfaehrt oben Druckspannungen
und unten Zugspannungen. Waehrend seiner Einsatzzeit koennte die
Lastgeschichte diese beiden Bereiche unterschiedlich schaedigen: Mikrorisswachstum,
Kriechen, viskoelastische Verformung und Ermuedunseffekte wirken sich in
Druck- und Zugzone unterschiedlich aus.

\section{Forschungsfrage}
\label{sec:research-question}

Die zentrale Forschungsfrage dieser Bachelorarbeit lautet:

\begin{quote}
\textit{Unterscheidet sich die Restfestigkeit zwischen der frueheren Druck-
und Zugzone eines gebrauchten Holzbalkens nach Jahrzehnten Belastungsgeschichte
signifikant?}
\end{quote}

Daraus ergeben sich spezifischere Fragen:

\begin{enumerate}
  \item Welche Materialphysik erklaert moegliche Unterschiede zwischen den
    Zonen?
  \item Welche Forschungsergebnisse existieren bereits zu Altholzfestigkeit
    und Zonenvariationen?
  \item Wie laesst sich eine statistisch aussagekraeftige Versuchsmethodik
    entwerfen?
\end{enumerate}

\section{Ziele und Beitrag}
\label{sec:contribution}

Diese Arbeit verfolgt folgende Ziele:

\begin{enumerate}
  \item Umfassende Recherche der aktuellen Rechtslage und Forschung zu
    Altholzfestigkeit
  \item Verstaendnis der Material- und Versagensphysik in Druck- vs. Zugzonen
  \item Entwurf einer Versuchsmethodik zum Festigkeitsvergleich nach EN- und
    DIN-Standards
  \item Preliminary-Ergebnisse zur Frage, ob ein signifikanter Unterschied
    vorhanden ist
\end{enumerate}

\section{Struktur der Arbeit}
\label{sec:structure}

Kapitel~\ref{ch:background} behandelt die Regulierung von Altholz und
Wiederverwendungsanforderungen (AltholzV, Klassifizierung). Kapitel~\ref{ch:literature}
fasst den Forschungsstand zusammen. Kapitel~\ref{ch:data} behandelt die
Materialphysik der Lastgeschichte. Kapitel~\ref{ch:strategy} entwirft die
Versuchsmethodik. Kapitel~\ref{ch:results} zeigt die Struktur der Ergebnisse,
und Kapitel~\ref{ch:discussion} bewertet sie kritisch.
